\documentclass[a4paper]{article}
\usepackage{amsfonts}
\usepackage{amsmath}
\usepackage{amssymb}
\usepackage{graphicx}  
\usepackage{cancel}

\newcommand{\degree}{^{\circ}}
\newcommand{\cis}{\operatorname{cis}}
\newcommand{\Bgsin}{\operatorname{Bgsin}}
\newcommand{\Bgcos}{\operatorname{Bgcos}}
\newcommand{\Bgtan}{\operatorname{Bgtan}}
\newcommand{\limit}{\lim\limits}
\newcommand{\summation}{\displaystyle\sum}
\newcommand{\dom}{\operatorname{dom}}
\newcommand{\Int}{\displaystyle\int}

\begin{document}
  
\noindent \large Auteur: Rune Suy \\
\noindent \large Studierichting: Informatica\\
\noindent \large Oefeningenreeks 6B nr. 4.2\\

\medskip

\normalsize

$x1=?$\\

geg: $x_n = -1", s_n=-45, v=-2$\\

$x_n = x_1 + (n-1) \cdot v$\\

$s_n = \dfrac{x_1 + x_n}{2} \cdot n$\\

$\Leftrightarrow \left\{
\begin{array}{l}
    -13 = x_1 + (n-1)\cdot(-2)\\
    -45 = \dfrac{x_1 - 13}{2} \cdot n
\end{array}
\right.$

$\Leftrightarrow \dfrac{-90}{x_1 -13}=n$ als $x_1 \neq 0$\\

$\Leftrightarrow -13 = x_1 + \dfrac{180 + 2(x_1 - 13)}{x_1 - 13}$\\

$\Leftrightarrow -13 = \dfrac{x_1^2 - 13x_1 + 180 + 2x_1 - 26}{x_1 - 13}$\\

$\Leftrightarrow -13x_1 + 169 = x_1^2 - 11x_1 + 154$\\

$\Leftrightarrow x_1^2 + 2x_1 - 1 = 0$\\

$D=64$\\

$\Leftarrow x_1 = -5$ of $x_1 = 3$\\

\begin{thebibliography}{999}
%\bibitem{a} Referentie
\end{thebibliography}
\end{document} 
